\sectionkicker{Theme synthesis}
\subsection*{横断比較 — 単発Launchでは追えないfamily lifecycle}
\addcontentsline{toc}{subsection}{横断比較 — 単発Launchでは追えないfamily lifecycle}
GPT-5系列は一回の発表で閉じず、base release、developer API、安全手法、後続世代、coding specializationへ分岐した。半年の変化を読むには、同じfamily名に属していてもevent identityを分離する必要がある。
\medskip
\begin{center}
\small
\renewcommand{\arraystretch}{1.16}
\begin{tabularx}{\linewidth}{@{}p{0.20\linewidth}X@{}}
\toprule
観察軸 & 一次資料・論文から読めること \\
\midrule
release packageの起点 & 8月のGPT-5はmodel release、developer API、safe-completionsなどを同日のrelease packageとして持つ一方、それぞれの事実と評価claimは別の境界で扱う。 \cite{src-aec54e92a7f3,src-05738f6c55b7,src-964121647555} \\
family更新 & 11月のGPT-5.1と12月のGPT-5.2は8月launchの追記ではなく、別の日付を持つmodel-family eventとして残す。family名の連続性とrelease identityは同じではない。 \cite{src-5bea303cfdb1,src-55da04f42e32} \\
coding specialization & Codex-MaxとGPT-5.2-Codexはfamily更新をcoding用途へspecializeした別eventとして位置付ける。base model releaseとcoding product/modelの公開を同じ出来事にはしない。 \cite{src-10ab1c39be14,src-cd8d6c8023f2} \\
評価claimの境界 & 各releaseで提示されたbenchmarkや優位性表現はvendor-attributed claimのまま維持する。一次資料から確定できるrelease、API availability、model identity、control情報とは分けて読む。 \cite{src-aec54e92a7f3,src-05738f6c55b7,src-964121647555,src-55da04f42e32} \\
\bottomrule
\end{tabularx}
\renewcommand{\arraystretch}{1.0}
\end{center}
\normalsize
\begin{claimboundary}[この比較の境界]
この表は本号で既に検証・参照している一次資料と論文を横断比較の形へ再配置したものである。vendor / project / author が報告した性能値や優位性は独立再現済みの一般則へ変換せず、本文とSource Notesの帰属境界をそのまま維持する。
\end{claimboundary}
